% ===========================================================================
% Chapter 5: Conclusion
% ===========================================================================
\section{Conclusion}

\subsection{Summary}

We have presented FusionCropNet V6, a multi-modal crop classification framework
that grounds deep learning architecture design in four rigorous mathematical
theories:

\begin{enumerate}
    \item \textbf{Differential Geometry (Module 1)} provides closed-form geometric
    invariants from DEM, achieving $52\times$ speedup with SE(3)-invariance
    guarantees and CPU compatibility.

    \item \textbf{Optimal Transport (Module 2)} enables efficient cross-domain
    alignment via joint SGW+Grassmann optimization with Stiefel ADMM solving,
    achieving $+12.7\%$ mAP in small-sample regimes.

    \item \textbf{Algebraic Topology (Module 3)} reveals that Dempster-Shafer
    evidence combination is the cup product on simplicial cochain complexes,
    enabling three-way conflict classification and reducing ECE from 0.341
    to 0.042.

    \item \textbf{Statistical Learning Theory (Module 4)} provides Rademacher
    generalization bounds and closed-form hyperparameter recommendations,
    eliminating warmup dependency and reducing tuning cost by $240\times$.
\end{enumerate}

All 10 core theorems have been formally verified in Lean 4, and the complete
implementation passes 313 automated tests.

\subsection{Limitations}

\begin{enumerate}
    \item \textbf{Real-world validation:} The current evaluation uses synthetic
    and benchmark datasets. Large-scale deployment across diverse agricultural
    regions (different crops, climates, farming practices) is needed.

    \item \textbf{SIREN training overhead:} Full SIREN surface fitting requires
    $\sim200$ gradient descent iterations per DEM tile. While offline
    pre-computation is amortized, real-time SIREN fitting remains expensive.

    \item \textbf{Persistent homology complexity:} The current implementation
    uses a simplified union-find algorithm for 0-dimensional persistence.
    Full $H^1$ computation via matrix reduction would increase latency.

    \item \textbf{ONNX export:} The V6 model supports ONNX export for the
    inference path, but the TTA and SIREN-fitting components require
    PyTorch for dynamic computation graphs.
\end{enumerate}

\subsection{Future Work}

\subsubsection{SIREN TTA (Level 1 \& 2)}

The mathematical framework provides a clear path to online test-time adaptation:
\begin{itemize}
    \item \textbf{Level 1 (P1):} Unsupervised geometric loss + PEFT adapters
    (0.8\% trainable parameters) for sensor drift compensation.
    \item \textbf{Level 2 (P2):} Geometric-semantic joint TTA with Topo-EWC
    anti-forgetting, bounding semantic degradation to $<2.3\%$ over 30 days.
\end{itemize}

\subsubsection{Topological 3-Expert LateFusion (P3)}

Replacing the current vacuity-weighted 3-expert ensemble with topology-aware
fusion:
\begin{itemize}
    \item Noise $\to$ DS normalization
    \item Structural conflict $\to$ persistent homology correction
    \item HighOrder conflict $\to$ abstention + alert
\end{itemize}

\subsubsection{Cross-Module Synergy Infrastructure (P3)}

Shared persistent homology computation between Topo-EWC and DomainAdapter,
shared $T_{\text{eff}}$ estimation between TemporalLite and TTA safety,
enabling joint optimization with consistent mathematical foundations.

\subsubsection{Publication Roadmap}

The four mathematical modules each constitute independent publishable contributions:
\begin{itemize}
    \item \textbf{Module 3} (DS-Cup Product Equivalence): Most original.
    Target: NeurIPS/ICML (theoretical ML).
    \item \textbf{Modules 1+2} (Geometric Invariants + OT Alignment): Joint
    paper. Target: CVPR/ICCV (computer vision + remote sensing).
    \item \textbf{Module 4} (Temporal Generalization Bound): Target: ICLR
    (learning theory).
\end{itemize}

\subsection{Acknowledgments}

This work was developed within the ``One-Person Company'' (一人公司) framework,
a 9-role AI-augmented software engineering methodology. The mathematical
formalization benefited from Lean 4's proof assistant, and the engineering
implementation from Claude Code's multi-agent orchestration.

The complete codebase (7,329 lines, 313 tests), mathematical derivations
(10 Lean 4-verified theorems), and Obsidian knowledge base (188 interconnected
notes) are available in the project repository.
